\section{Experiment 2 - Nanoscale Optical Metrology with SMLM}

In this experiment, you will image DNA origami nanorulers as part of a metrology experiment. The goal is to asses the precision of the microscope and the SMLM technique.

\subsection{Materials}

\begin{itemize}
    \item{DNA origami nanoruler sample}
    \item{Immersion oil}
    \item{Lens cleaning products (see \autoref{sec:remove-oil})}
\end{itemize}

\subsection{Description of the Sample}

The sample is a DNA origami nanoruler and is illustrated in \autoref{fig:dna-origami-nanoruler}. It consists of a long DNA strand that has been folded into a specific shape. The shape is designed to be imaged with SMLM. The nanoruler is labeled with fluorescent markers at specific locations. The distance between the markers is known to a high degree of precision. The nanoruler is mounted on a glass coverslip and imaged as in \autoref{sec:mount-sample}.

\begin{figure}
    \centering
    \includegraphics[width=0.75\textwidth]{dna-origami-nanoruler.png}
    \caption{Illustration of the DNA origami nanoruler. The nanoruler has two fluorescent markers at a known separation distance $L$.}
    \label{fig:dna-origami-nanoruler}
\end{figure}

In this laboratory, the separation between fluorophores is 94 nm.

\subsection{Procedure}

\textbf{TO DO}
